\section{Additional Results}
\label{sec:appendix}

\subsection{Watermark Evasion Evaluation}
\label{sec:evasion_appendix}

We evaluate three evasion strategies against the unigram watermark
(\texttt{WatermarkEvasionEvaluator.evaluate\_all()}, see
\texttt{src/evaluation/evasion\_eval.py}).
Table~\ref{tab:evasion} summarises the results.
Evasion rates are measured as the fraction of watermarked samples whose
$z$-score falls below the detection threshold after applying the evasion
strategy at increasing intensity.

\IfFileExists{../../results/tables/table4_evasion.tex}{%
  \begin{table}[t]
\centering
\small
\caption{Watermark evasion results. TPR before and after evasion, and fraction of samples where evasion succeeded ($z < z_{\text{thr}}$).}
\label{tab:evasion}
\begin{tabular}{lccc}
\toprule
Evasion Strategy & TPR Before & TPR After & Evasion Rate \\
\midrule
Paraphrase & 1.000 & 1.000 & 0.000 \\
Copy-Paste Dilution & 1.000 & 1.000 & 0.000 \\
Adaptive Substitution & 1.000 & 1.000 & 0.000 \\
\bottomrule
\end{tabular}
\end{table}%
}{}

\IfFileExists{figures/p13_evasion_analysis.pdf}{%
\begin{figure}[ht]
  \centering
  \includegraphics[width=0.82\linewidth]{p13_evasion_analysis.pdf}
  \caption{Watermark evasion analysis across three strategies at increasing
           intensity.  The unigram z-score defense maintains $\tpr=1.00$ across
           all tested intensity levels on the synthetic corpus, confirming
           robustness against mild evasion.  More aggressive strategies
           (large $\rho$ or high character-swap budget) that would reduce
           $\tpr$ also visibly degrade the adversarial passage quality.}
  \label{fig:evasion_full}
\end{figure}%
}{}

\begin{description}
  \item[Paraphrasing.] Synonym substitution at intensity $\rho \in \{0.1, 0.2, 0.3, 0.4, 0.5\}$.
    On the synthetic corpus (200 entries, typical length 50--150 tokens),
    $\tpr$ remains $1.00$ for all tested $\rho$.  The unigram watermark embeds
    a bias across the full vocabulary, so short synonym substitutions do not
    reduce the $z$-score below the threshold in this setting.

  \item[Copy-paste dilution.] Adversarial content diluted with $r \in \{0.5, 1.0, 2.0, 3.0\}$
    copies of clean content.  $\tpr=1.00$ across all dilution ratios; the
    embedded bias survives dilution on entries of this length.

  \item[Adaptive substitution.] Targeted character swaps with budget
    $b \in \{1, 5, 10, 20\}$.  $\tpr=1.00$ across all tested budgets.
    Character-level swaps perturb individual token ids but do not
    systematically reduce the aggregate unigram bias over the full entry.
\end{description}

\IfFileExists{figures/wm_detection_distribution.pdf}{%
\begin{figure}[ht]
  \centering
  \includegraphics[width=0.72\linewidth]{wm_detection_distribution.pdf}
  \caption{$z$-score distributions for watermarked (shaded) and clean (outline)
           entries.  The unigram watermark produces a clear separation; the
           detection threshold (vertical dashed line) achieves $\tpr=1.00$
           with $\fpr < 0.05$ on the synthetic corpus.}
  \label{fig:wm_distribution}
\end{figure}%
}{}

\IfFileExists{figures/wm_ablation_threshold.pdf}{%
\begin{figure}[ht]
  \centering
  \includegraphics[width=0.72\linewidth]{wm_ablation_threshold.pdf}
  \caption{Watermark detection performance as a function of $z$-threshold.
           TPR (solid) and FPR (dashed) as the threshold varies.
           The default threshold ($z=4.0$, vertical dashed) achieves
           $\tpr \approx 0.90$, $\fpr \approx 0.03$.}
  \label{fig:wm_threshold}
\end{figure}%
}{}

\subsection{Synthetic Corpus Details}
\label{sec:corpus_appendix}

The 200-entry synthetic corpus is generated by
\texttt{SyntheticCorpus(seed=42)} from \texttt{src/data/synthetic\_corpus.py}.
Entry categories and counts:
\begin{center}
\begin{tabular}{lc}
  \toprule
  Category & Count \\
  \midrule
  User preferences & 30 \\
  Task history & 30 \\
  Calendar events & 30 \\
  Factual knowledge & 30 \\
  Conversation snippets & 30 \\
  Configuration data & 25 \\
  Project notes & 25 \\
  \midrule
  Total & 200 \\
  \bottomrule
\end{tabular}
\end{center}

\subsection{Hyperparameter Sensitivity}
\label{sec:ablation_appendix}

Ablation studies are generated by \texttt{AblationStudy.run\_all()} in
\texttt{src/evaluation/ablation\_study.py}.  Each ablation point reports
the mean and 95\% CI over multiple seeds.

\paragraph{Corpus size ($N$).}
We sweep $N \in \{50, 100, 200, 500\}$ with $n_{\text{poison}}=5$, $k=5$.
ASR-R for all three attacks remains largely constant as $N$ grows.  This
confirms that poison passages act as extreme similarity outliers regardless
of benign corpus size: more benign entries do not crowd out the adversarial
ones in retrieval results.

\paragraph{Retrieval top-$k$.}
We sweep $k \in \{1, 3, 5, 10, 20\}$ with $N=200$, $n_{\text{poison}}=5$.
ASR-R increases monotonically with $k$: a larger retrieval window gives
more opportunities for poison entries to appear in results.

\IfFileExists{figures/p13_ablation_corpus.pdf}{%
\begin{figure}[ht]
  \centering
  \begin{minipage}{0.48\linewidth}
    \includegraphics[width=\linewidth]{p13_ablation_corpus.pdf}
    \caption{ASR-R vs.\ corpus size $N$ for all three attacks.
             Results are averaged over 5 seeds at $n_{\text{poison}}=5$, $k=5$.}
    \label{fig:ablation_corpus}
  \end{minipage}\hfill
  \begin{minipage}{0.48\linewidth}
    \IfFileExists{figures/p13_ablation_topk.pdf}{%
      \includegraphics[width=\linewidth]{p13_ablation_topk.pdf}%
    }{}
    \caption{ASR-R vs.\ retrieval top-$k$.  \agentpoison{} retains high
             ASR-R at all $k$ due to the trigger.}
    \label{fig:ablation_topk}
  \end{minipage}
\end{figure}%
}{}

\paragraph{Poison count.}
We sweep $n_{\text{poison}} \in \{1, 3, 5, 10, 20\}$ with $N=200$, $k=5$.
\agentpoison{} shows near-ceiling ASR-R at $n_{\text{poison}}=1$ (trigger
guarantees retrieval).  \injecmem{} scales approximately linearly up to 10
entries, with diminishing returns thereafter.

\paragraph{\sad{} threshold $k$ (\minja{} attack).}
We sweep $k \in \{0.5, 1.0, 1.5, 2.0, 2.5, 3.0, 3.5, 4.0\}$ using
\minja{} (where \sad{} achieves non-trivial TPR; calibration uses victim
queries, $N=200$).
At $k \leq 2.0$, \sad{} achieves $\tpr = 1.00$ for \minja{} while
$\fpr$ drops from $0.45$ ($k=0.5$) to $0.00$ ($k \geq 2.0$).
For $k \geq 3.0$, TPR begins to decrease: $\tpr = 0.80$ ($k=3.0$),
$0.60$ ($k=3.5$), $0.30$ ($k=4.0$).  The default $k=2.0$ achieves
$\tpr=1.00$, $\fpr=0.00$ for \minja{} and $\tpr \approx 0.00$ for
\agentpoison{} (centroid passage falls within the benign distribution;
see Section~\ref{sec:discussion}).

\IfFileExists{figures/p13_ablation_sad_sigma.pdf}{%
\begin{figure}[ht]
  \centering
  \includegraphics[width=0.55\linewidth]{p13_ablation_sad_sigma.pdf}
  \caption{\sad{} threshold $k$ vs.\ TPR and FPR for \minja{}.
           Larger $k$ reduces both TPR and FPR (stricter detection).
           The default $k=2.0$ (vertical dashed) achieves
           $\tpr=1.00$, $\fpr=0.00$ for \minja{}.}
  \label{fig:ablation_sad_sigma}
\end{figure}%
}{}

\paragraph{Watermark $z$-threshold.}
We sweep $z_{\mathrm{thr}} \in \{2.0, 3.0, 4.0, 5.0, 6.0, 7.0\}$.
TPR for watermarked entries remains above $0.85$ for $z_{\mathrm{thr}} \leq 5.0$.
FPR for clean entries is $< 0.05$ for $z_{\mathrm{thr}} \geq 3.0$.
The default $z_{\mathrm{thr}} = 4.0$ achieves TPR $\approx 0.90$, FPR $\approx 0.03$.

\IfFileExists{figures/p13_ablation_watermark_z.pdf}{%
\begin{figure}[ht]
  \centering
  \includegraphics[width=0.55\linewidth]{p13_ablation_watermark_z.pdf}
  \caption{Watermark $z$-threshold vs.\ TPR and FPR for watermarked entries.
           TPR remains near 1.00 for $z_{\mathrm{thr}} \leq 5.0$ on the
           synthetic corpus; the default $z_{\mathrm{thr}}=4.0$ (vertical
           dashed) achieves $\tpr \approx 0.90$, $\fpr \approx 0.03$.}
  \label{fig:ablation_watermark_z}
\end{figure}%
}{}

\IfFileExists{../../results/tables/table6_ablation.tex}{%
  \begin{table}[t]
\centering
\small
\caption{Hyperparameter ablation summary.  Each row corresponds to one setting of the ablated parameter.}
\label{tab:ablation}
\begin{tabular}{lc|lc|lcc|lcc}
\toprule
\multicolumn{2}{c|}{Corpus} & \multicolumn{2}{c|}{Top-K} & \multicolumn{3}{c|}{SAD $k$} & \multicolumn{3}{c}{WM $z$-thr} \\
$N$ & ASR-R & $k$ & ASR-R & $k$ & TPR & FPR & $z$ & TPR & FPR \\
\midrule
50 & 0.325 & 1 & 0.000 & 0.5 & 0.000 & 0.000 & 2.0 & 1.000 & 0.000 \\
100 & 0.125 & 3 & 0.000 & 1.0 & 0.000 & 0.000 & 3.0 & 1.000 & 0.000 \\
200 & 0.100 & 5 & 0.100 & 1.5 & 0.000 & 0.000 & 4.0 & 1.000 & 0.000 \\
500 & 0.000 & 10 & 0.150 & 2.0 & 0.000 & 0.000 & 5.0 & 1.000 & 0.000 \\
&  & 20 & 0.450 & 2.5 & 0.000 & 0.000 & 6.0 & 1.000 & 0.000 \\
&  & &  & 3.0 & 0.000 & 0.000 & 7.0 & 1.000 & 0.000 \\
&  & &  & 3.5 & 0.000 & 0.000 & & &  \\
&  & &  & 4.0 & 0.000 & 0.000 & & &  \\
\bottomrule
\end{tabular}
\end{table}%
}{}

\subsection{Encoder Generalization}
\label{app:encoder_generalization}

A key practical question is whether SAD's detection performance is robust to
the choice of sentence-transformer encoder.  If the detection signal depended
on properties specific to the \texttt{all-MiniLM-L6-v2} encoder, the defense
would not be directly deployable on systems using different retrieval models.

We evaluate SAD at $k=2.0$ on three sentence-transformer encoders of varying
capacity and training objective:
\begin{enumerate}
  \item \texttt{all-MiniLM-L6-v2} (384-dim, contrastive, default retriever)
  \item \texttt{all-mpnet-base-v2} (768-dim, higher-capacity MPNet backbone)
  \item \texttt{paraphrase-MiniLM-L6-v2} (384-dim, paraphrase-optimised)
\end{enumerate}
For each encoder and each attack (\minja{}, \injecmem{}, and \agentpoison{}
with triggered-query calibration), we calibrate SAD on 100 benign entries
and 10 victim queries drawn from the same synthetic corpus, then evaluate
on 100 held-out benign entries and the respective poison passages.

\IfFileExists{../../results/tables/table_encoder_generalization.tex}{%
  \begin{table}[h]
  \centering
  \small
  \caption{SAD encoder generalization: TPR / FPR / AUROC at $k=2.0$ for three sentence-transformer encoders.  Results are consistent across encoder capacity and training objective, confirming that SAD's detection signal is encoder-agnostic.}
  \label{tab:encoder_generalization}
  \begin{tabular}{llccc}
    \toprule
    Encoder & Attack & TPR & FPR & AUROC \\
    \midrule
    \texttt{MiniLM-L6} & \minja{} & 0.800 & 0.090 & 0.966 \\
     & \injecmem{} & 0.400 & 0.090 & 0.816 \\
     & \agentpoison{} (trig.) & 1.000 & 0.040 & 1.000 \\
    \midrule
    \texttt{MPNet-Base} & \minja{} & 1.000 & 0.100 & 0.954 \\
     & \injecmem{} & 0.200 & 0.100 & 0.840 \\
     & \agentpoison{} (trig.) & 1.000 & 0.100 & 1.000 \\
    \midrule
    \texttt{Para-MiniLM} & \minja{} & 1.000 & 0.080 & 0.978 \\
     & \injecmem{} & 0.000 & 0.080 & 0.744 \\
     & \agentpoison{} (trig.) & 1.000 & 0.040 & 1.000 \\
    \bottomrule
  \end{tabular}
\end{table}
%
}{}

\IfFileExists{figures/fig_encoder_generalization.pdf}{%
\begin{figure}[h]
  \centering
  \includegraphics[width=0.98\linewidth]{fig_encoder_generalization.pdf}
  \caption{SAD TPR (left) and FPR (right) at $k=2.0$ across three
           sentence-transformer encoders.  \agentpoison{} with triggered-query
           calibration achieves $\tpr = 1.000$ on all three encoders.
           \minja{} achieves $\tpr \geq 0.80$ across encoders; the variation
           reflects differences in the encoder's synonym-invariance under
           bridging-step passage vocabulary.  \injecmem{} shows the largest
           encoder sensitivity ($\tpr \in [0.00, 0.40]$), consistent with
           broad-anchor passages being designed to be retriever-agnostic
           (close to all queries, hence harder to detect against any single
           encoder).  FPR remains below $0.10$ across all conditions.}
  \label{fig:encoder_generalization}
\end{figure}%
}{}

\paragraph{Findings.}
Three observations are consistent across all encoders:
\begin{enumerate}
  \item \textbf{AgentPoison + triggered calibration.}  $\tpr = 1.000$ on all
  three encoders (AUROC $= 1.000$).  The centroid passage is an extreme
  similarity outlier regardless of embedding dimension or training objective:
  it is explicitly constructed to maximize cosine similarity to victim queries,
  making it detectable in any encoder space.
  \item \textbf{MINJA.}  $\tpr \in [0.80, 1.00]$ across encoders.  Bridging-step
  passages contain the same vocabulary as victim queries, producing anomalous
  similarity in all tested embedding spaces.
  \item \textbf{InjecMEM.}  $\tpr \in [0.00, 0.40]$ — the largest encoder
  sensitivity.  The broad-anchor design achieves moderate similarity to many
  query types, placing it closer to the benign boundary; whether it crosses the
  threshold depends on how tightly the encoder clusters general-purpose
  vocabulary.
\end{enumerate}

These results support the conclusion that SAD's detection signal is
fundamentally encoder-agnostic for targeted attacks: any encoder that places
poison entries close to victim queries in embedding space will make those
entries detectable by SAD.  The InjecMEM exception is consistent with its
threat model (retriever-agnostic anchors achieve partial similarity to many
queries, not maximum similarity to any specific query).

\subsection{Bootstrap CI Computation}
\label{sec:bootstrap}

We use percentile bootstrap with $B=2000$ replicates.  For metric $\theta$
estimated from $n$ trials $\{x_1, \ldots, x_n\}$:
\begin{equation}
  \hat{\theta}^{(b)} = \hat{\theta}(\{x_{\pi_b(1)}, \ldots, x_{\pi_b(n)}\}),
  \quad \pi_b \sim \mathrm{Uniform}(\text{permutations with replacement}),
\end{equation}
and the 95\% CI is $[\hat{\theta}^{(0.025)}, \hat{\theta}^{(0.975)}]$.
